\par\SourceLeaf{222}
Sed quoties $1-(\gamma-k)$ sive $k+1-\gamma$ est quantitas positiva, demonstravimus
esse (formula 85)
\[
f(\alpha-k,\beta-k,\gamma-k)
=\frac{\Pi(\alpha+\beta-\gamma-k)\Pi(k-\gamma)}{\Pi(\alpha-\gamma)\Pi(\beta-\gamma)}
\]
Quare quum $k$, quidquid sit $\gamma$, semper accipi possit tantus, ut $k+1-\gamma$ evadat
quantitas positiva, erit \textit{generaliter}
\[
f(\alpha,\beta,\gamma)
=\frac{\Pi(\alpha+\beta-\gamma)\Pi(-\gamma)}{\Pi(\alpha-\gamma)\Pi(\beta-\gamma)}
\]
et proin
\[
f(\alpha+1-\gamma,\beta+1-\gamma,2-\gamma)
=\frac{\Pi(\alpha+\beta-\gamma)\Pi(\gamma-2)}{\Pi(\alpha-1)\Pi(\beta-1)}
\]
ita ut formula nostra fiat
\[
\begin{aligned}[c]
[86]\quad &F(\alpha,\beta,\alpha+\beta+1-\gamma,1-x)\\
&=\frac{\Pi(\alpha+\beta-\gamma)\Pi(-\gamma)}{\Pi(\alpha-\gamma)\Pi(\beta-\gamma)}F(\alpha,\beta,\gamma,x)\\
&\quad+\frac{\Pi(\alpha+\beta-\gamma)\Pi(\gamma-2)}{\Pi(\alpha-1)\Pi(\beta-1)}x^{1-\gamma}\\&\qquad(1-x)^{\gamma-\alpha-\beta}F(1-\alpha,1-\beta,2-\gamma,x)
\end{aligned}
\]
sive mutata $\gamma$ in $\alpha+\beta+1-\gamma$
\[
\begin{aligned}[c]
[87]\quad &F(\alpha,\beta,\gamma,1-x)\\
&=\frac{\Pi(\gamma-1)\Pi(\gamma-\alpha-\beta-1)}{\Pi(\gamma-\alpha-1)\Pi(\gamma-\beta-1)}F(\alpha,\beta,\alpha+\beta+1-\gamma,x)\\
&\quad+\frac{\Pi(\gamma-1)\Pi(\alpha+\beta-\gamma-1)}{\Pi(\alpha-1)\Pi(\beta-1)}x^{\gamma-\alpha-\beta}\\&\qquad(1-x)^{1-\gamma}F(1-\alpha,1-\beta,\gamma+1-\alpha-\beta,x)
\end{aligned}
\]

Si magis placet, scribere licet

in formula 86
\[
\begin{aligned}&\text{pro }(1-x)^{\gamma-\alpha-\beta}F(1-\alpha,1-\beta,2-\gamma,x)\\&\ .\ .\ .\ .\ F(\alpha+1-\gamma,\beta+1-\gamma,2-\gamma,x)\end{aligned}
\]
in formula 87
\[
\begin{aligned}&\text{pro }(1-x)^{1-\gamma}F(1-\alpha,1-\beta,\gamma+1-\alpha-\beta,x)\\&\ .\ .\ .\ .\ F(\gamma-\alpha,\gamma-\beta,\gamma+1-\alpha-\beta,x)\end{aligned}
\]

\SixthArticle{44}

Quoties itaque elemento quarto in aliqua. serie sub forma nostra contenta
valor tribuitur inter $0{,}5$ et $1$, convergentiae lentiori per formulas praecedentes
remedium affertur, quippe quae illam in duas alias series similes dispescunt,
\SourceLeaf{223}
eo citius convergentes, quo tardius illa convergebat. Sed excipere oportet casus
speciales, ubi haec transformatio non succedit, quoties scilicet in serie transformanda differentia inter elementum tertium summamque duorum elementorum primorum fit numerus integer. Si enim in formula 86 $\gamma$ est $=0$ vel aequalis numero negativo integro, manifesto $F(\alpha,\beta,\gamma,x)$ fit series inepta (art. 2)
atque factor $\Pi(\gamma-2)$ infinitus; si vero $\gamma$ est integer positivus unitate major
$F(1-\alpha,1-\beta,2-\gamma,x)$ atque $F(\alpha+1-\gamma,\beta+1-\gamma,2-\gamma,x)$ fiunt series ineptae et $\Pi(-\gamma)$ infinitus; denique si $\gamma=1$ duae series transformatae $F(\alpha,\beta,\gamma,x)$
atque $F(1-\alpha,1-\beta,2-\gamma,x)$ vel $F(\alpha+1-\gamma,\beta+1-\gamma,2-\gamma,x)$, quae ideo
cum $F(\alpha,\beta,\gamma,x)$ identica evadit, hocce quidem incommodo non laborant, sed
nihilominus transformatio nullius est usus, quum utraque series transformata per
co\"efficientem infinitum $\Pi(-1)$ multiplicata sit. Operae itaque pretium erit ostendere, quomodo in his quoque casibus convergentia lentior in citiorem mutari possit.

\SixthArticle{45}

Sit $k$ numerus integer positivus (sive etiam $=0$) designemusque $k+1$
primos terminos seriei $F(\alpha,\beta,\gamma,x)$ per $X$. Terminus sequens erit
\[
=\frac{\alpha.\alpha+1.\alpha+2\ .\ .\ .\ .\ \alpha+k.\beta.\beta+1.\beta+2\ .\ .\ .\ .\ \beta+k}
{1\ .\ 2\ .\ 3\ .\ .\ .\ .\ k+1.\gamma.\gamma+1.\gamma+2\ .\ .\ .\ .\ \gamma+k}x^{k+1}
\]
qui etiam ita exhiberi potest
\[
\frac{\Pi(\gamma-1)}{\Pi(\alpha-1)\Pi(\beta-1)}\ .\
\frac{\Pi(\alpha+k)\Pi(\beta+k)}{\Pi(k+1)\Pi(\gamma+k)}x^{k+1}
\]
similique modo termini sequentes. Hinc colligitur
\[
\mathrm{I.}\qquad
\frac{\Pi(\alpha+\beta-\gamma)(\Pi(-\gamma)}{\Pi(\alpha-\gamma)\Pi(\beta-\gamma)}\ .\ F(\alpha,\beta,\gamma,x)
\quad\text{exprimi posse per}
\]
\[
\begin{aligned}&\frac{\Pi(\alpha+\beta-\gamma)\Pi(-\gamma)}{\Pi(\alpha-\gamma)\Pi(\beta-\gamma)}\ .\ X\\&+\frac{\Pi(\alpha+\beta-\gamma)\Pi(-\gamma)\Pi(\gamma-1)}{\Pi(\alpha-1)\Pi(\beta-1)\Pi(\alpha-\gamma)\Pi(\beta-\gamma)}\\&\qquad\Sigma\left\{\frac{\Pi(\alpha+k+t)\Pi(\beta+k+t)}{\Pi(k+t+1)\Pi(\gamma+k+t)}x^{k+1+t}\right\}\end{aligned}
\]
si pro $t$ omnes valores $0,1,2,3$ etc in infinitum substitui concipiuntur.

Simili modo $F(\alpha+1-\gamma,\beta+1-\gamma,2-\gamma,x)$ exprimi potest per
\[
\frac{\Pi(1-\gamma)}{\Pi(\alpha-\gamma)\Pi(\beta-\gamma)}
\Sigma\left(\frac{\Pi(\alpha-\gamma+t)\Pi(\beta-\gamma+t)}{\Pi t\,\Pi(1-\gamma+t)}x^t\right)
\]
$t$ perinde ut ante determinato, adeoque quum sit $\Pi(1-\gamma)=(1-\gamma)\Pi(-\gamma)$ atque $\Pi(\gamma-1)=-(1-\gamma)\Pi(\gamma-2)$, patet esse
\SourceLeaf{224}
\[
\begin{aligned}[c]
\mathrm{II.}\qquad&\frac{\Pi(\alpha+\beta-\gamma)\Pi(\gamma-2)}{\Pi(\alpha-1)\Pi(\beta-1)}x^{1-\gamma}F(\alpha+1-\gamma,\beta+1-\gamma,2-\gamma,x)\\
&=-\frac{\Pi(\alpha+\beta-\gamma)\Pi(-\gamma)\Pi(\gamma-1)}{\Pi(\alpha-1)\Pi(\beta-1)\Pi(\alpha-\gamma)\Pi(\beta-\gamma)}
\\&\qquad\Sigma\left\{\frac{\Pi(\alpha-\gamma+t)\Pi(\beta-\gamma+t)}{\Pi t\,\Pi(1-\gamma+t)}x^{1+t-\gamma}\right\}
\end{aligned}
\]
Hinc formula 86 etiam ita exhiberi potest:
\[
\begin{aligned}[c]
&F(\alpha,\beta,\alpha+\beta+1-\gamma,1-x)\\
&=\frac{\Pi(\alpha+\beta-\gamma)\Pi(-\gamma)}{\Pi(\alpha-\gamma)\Pi(\beta-\gamma)}X\\
&\quad+\frac{\Pi(\alpha+\beta-\gamma)\Pi(-\gamma)\Pi(\gamma-1)}{\Pi(\alpha-1)\Pi(\beta-1)\Pi(\alpha-\gamma)\Pi(\beta-\gamma)}
\\&\qquad\Sigma\Bigl\{\frac{\Pi(\alpha+k+t)\Pi(\beta+k+t)}{\Pi(k+t+1)\Pi(\gamma+k+t)}x^{k+1+t}
\\&\qquad-\frac{\Pi(\alpha-\gamma+t)\Pi(\beta-\gamma+t)}{\Pi(t-\gamma+1)\Pi t}x^{1+t-\gamma}\Bigr\}
\end{aligned}
\]
Haecce expressio protinus ostendit, singulas differentias, quae sunt sub signo $\Sigma$,
fieri $=0$, si supponatur $\gamma=-k$, sed quum hic simul fiat $\Pi(\gamma-1)$ quantitas infinite magna, productum finitum evadere posse patet. Cuius valorem ut
per quantitates finitas exprimamus, statuamus primo $\gamma+k=\omega$, unde fit
\[
\Pi(\gamma-1).\gamma.(\gamma+1)(\gamma+2)\ .\ .\ .\ .\ (\gamma+k-1)\omega=\Pi\omega
\]
sive
\[
\Pi(\gamma-1)=\frac{\Pi\omega}{\omega(\omega-1)(\omega-2)\ .\ .\ .\ .\ (\omega-k)}
\]
Rei summa vertitur itaque in eo, ut videamus, quid fiat
\[
\begin{aligned}&\frac{1}{\omega}\Bigl\{\frac{\Pi(\alpha-\gamma+t+\omega)\Pi(\beta-\gamma+t+\omega)}{\Pi(t-\gamma+1+\omega)\Pi(t+\omega)}x^{1+t-\gamma+\omega}\\&\qquad-\frac{\Pi(\alpha-\gamma+t)\Pi(\beta-\gamma+t)}{\Pi(t-\gamma+1)\Pi t}x^{1+t-\gamma}\Bigr\}\end{aligned}
\]
si $\omega$ in infinitum decrescat. Per principia nota autem hinc resultat
\[
-\frac{\mathrm{d}U}{\mathrm{d}\gamma}
\]
si brevitatis caussa statuimus
\[
\frac{\Pi(\alpha-\gamma+t)\Pi(\beta-\gamma+t)}{\Pi(t-\gamma+1)\Pi(t-k-\gamma)}x^{1+t-\gamma}=U
\]
solamque $\gamma$ tamquam variabilem spectamus. Sed hinc fit
\[
\begin{aligned}&\frac{\mathrm{d}U}{U\mathrm{d}\gamma}
=-\Psi(\alpha-\gamma+t)-\Psi(\beta-\gamma+t)\\&\qquad+\Psi(t-\gamma+1)+\Psi(t-k-\gamma)-\log x\end{aligned}
\]
Hinc colligitur pro $\gamma=-k$
\SourceLeaf{225}
\[
\begin{aligned}[c]
[88]\quad &F(\alpha,\beta,\alpha+\beta+1+k,1-x)\\
&=\frac{\Pi(\alpha+\beta+k)\Pi k}{\Pi(\alpha+k)\Pi(\beta+k)}X\\
&\quad+\frac{\Pi(\alpha+\beta+k)\Pi k}{\Pi(\alpha-1)\Pi(\beta-1)\Pi(\alpha+k)\Pi(\beta+k)(-1)(-2)\ldotp\ldotp\ldotp\ldotp(-k)}\\&\qquad\Sigma\{(\log x\\
&\qquad+\Psi(\alpha+t+k)+\Psi(\beta+t+k)-\Psi(t+k+1)-\Psi t)
\\&\qquad\frac{\Pi(\alpha+t+k)\Pi(\beta+t+k)}{\Pi(t+k+1)\Pi t}x^{1+t+k}\}\\
&=\frac{\Pi(\alpha+\beta+k)\Pi k}{\Pi(\alpha+k)\Pi(\beta+k)}X
\pm\frac{\Pi(\alpha+\beta+k)x^{1+k}}{\Pi(\alpha-1)\Pi(\beta-1)\Pi(k+1)}Y
\end{aligned}
\]
ubi
\[
\begin{aligned}[c]
&Y=\{\log x+\Psi(\alpha+k)+\Psi(\beta+k)-\Psi(k+1)-\Psi(0)\}\\&\qquad F(\alpha+k+1,\beta+k+1,k+2,x)\\
&+A\frac{\alpha+k+1.\beta+k+1}{1. k+2}x\\
&+(A+B)\frac{\alpha+k+1.\alpha+k+2.\beta+k+1.\beta+k+2}{1.2. k+2. k+3}xx\\
&+(A+B+C)\\&\qquad\frac{\alpha+k+1.\alpha+k+2.\alpha+k+3.\beta+k+1.\beta+k+2.\beta+k+3}{1.2.3. k+2. k+3. k+4}x^3\\
&+\text{etc.}
\end{aligned}
\]
atque
\[
\begin{aligned}[c]
A&=\frac{1}{\alpha+k+1}+\frac{1}{\beta+k+1}-\frac{1}{k+2}-1\\
B&=\frac{1}{\alpha+k+2}+\frac{1}{\beta+k+2}-\frac{1}{k+3}-\frac12\\
C&=\frac{1}{\alpha+k+3}+\frac{1}{\beta+k+3}-\frac{1}{k+4}-\frac13,\\
&\text{etc.}
\end{aligned}
\]
signumque superius vel inferius accipiendum est, prout $k$ est numerus par vel
impar.

\SixthArticle{46}
Hoc itaque modo $F(\alpha,\beta,\alpha+\beta+1-\gamma,1-x)$ transmutatur, si $\gamma$ est $0$
vel integer negativus. Casum $\gamma=+1$ prorsus simili modo tractare possumus,
sive brevius statuere possumus in ratiociniis praecedentidus $k=-1$, unde $X$
prorsus evanescit atque obtinemus:
\SourceLeaf{226}
\[
\begin{aligned}[c]
[89]\quad &F(\alpha,\beta,\alpha+\beta,1-x)\\
&=-\frac{\Pi(\alpha+\beta-1)}{\Pi(\alpha-1)\Pi(\beta-1)}
\\&\qquad\{\log x+\Psi(\alpha-1)+\Psi(\beta-1)-2\Psi(0)\}F(\alpha,\beta,1,x)\\
&\quad-\frac{\Pi(\alpha+\beta-1)}{\Pi(\alpha-1)\Pi(\beta-1)}
\{A\frac{\alpha.\beta}{1}x\\
&\quad+(A+B)\frac{\alpha.\alpha+1.\beta.\beta+1}{1.2.1.2}xx\\
&\quad+(A+B+C)\frac{\alpha.\alpha+1.\alpha+2.\beta.\beta+1.\beta+2}{1.2.3.1.2.3}x^3\\
&\quad+\text{etc.}\}
\end{aligned}
\]
ubi
\[
\begin{aligned}&A=\frac1\alpha+\frac1\beta-2,\\&
B=\frac1{\alpha+1}+\frac1{\beta+1}-\frac22.\\&
C=\frac1{\alpha+2}+\frac1{\beta+2}-\frac23,\\&\text{etc.}\end{aligned}
\]
Ita e.\ g.\ pro $\alpha=\frac12$, $\beta=\frac12$ obtinemus (cfr.\ form.\ 52, 71)
\[
\begin{aligned}[c]
[90]\quad &F(\tfrac12,\tfrac12,1,1-x)\\
&=-\frac1\pi\log\tfrac1{16}x. F(\tfrac12,\tfrac12,1,x)\\
&\quad-\frac1\pi\{2.\frac{1.1}{2.2}x
\\&\qquad+(2+\tfrac13)\frac{1.1.3.3}{2.2.4.4}xx
\\&\qquad+(2+\tfrac13+\tfrac2{15})\frac{1.1.3.3.5.5}{2.2.4.4.6.6}x^3\\
&\qquad+(2+\tfrac4{3.4}+\tfrac4{5.6}+\tfrac4{7.8})
\frac{1.1.3.3.5.5.7.7}{2.2.4.4.6.6.8.8}x^4+\text{etc.}\}\\
&=-\frac1\pi\{\log\tfrac1{16}xF(\tfrac12,\tfrac12,1,x)
\\&\qquad+\tfrac12x+\tfrac{21}{64}xx+\tfrac{185}{768}x^3+\tfrac{18655}{98304}x^4\\
&\qquad+\tfrac{102501}{655360}x^5+\tfrac{1394239}{10485760}x^6+\text{etc.}\}
\end{aligned}
\]
Denique casum tertium, ubi $\gamma$ est integer positivus unitate maior, seorsim
tractare haud necesse est, quum sit
\[
\begin{aligned}&F(\alpha,\beta,\alpha+\beta+1-\gamma,1-x)\\&=x^{\gamma-1}F(\alpha+1-\gamma,\beta+1-\gamma,\alpha+\beta+1-\gamma,1-x)\end{aligned}
\]
transformatioque seriei $F(\alpha+1-\gamma,\beta+1-\gamma,\alpha+\beta+1-\gamma,1-x)$ pro $\gamma>1$
ad casum primum sponte reducatur.

\SixthArticle{47}
Transimus ad alias transformationes, inter quas primum locum obtineat
substitutio $x=\frac{y}{y-1}$. Hinc fit $\mathrm{d}x=-\frac{\mathrm{d}y}{(y-1)^2}$, adeoque
\[
\frac{\mathrm{d}P}{\mathrm{d}x}=-\frac{\mathrm{d}P}{\mathrm{d}y}(1-y)^2,
\quad\text{differentiando denuo fit}
\]
\[
\mathrm{d}\frac{\mathrm{d}P}{\mathrm{d}x}=-(1-y)^2\mathrm{d}\frac{\mathrm{d}P}{\mathrm{d}y}
+2(1-y)\mathrm{d}P,\quad\text{adeoque}
\]
\[
\frac{\mathrm{dd}P}{\mathrm{d}x^2}=+(1-y)^4\frac{\mathrm{dd}P}{\mathrm{d}y^2}
-2(1-y)^3\frac{\mathrm{d}P}{\mathrm{d}y}
\]
\SourceLeaf{227}
Quibus valoribus substitutis transit aequatio 80 in hanc
\[
\begin{aligned}&0=\alpha\beta P+(1-y)(\gamma+(\alpha+\beta-1-\gamma)y)\frac{\mathrm{d}P}{\mathrm{d}y}\\&+(1-y)(y-yy)\frac{\mathrm{dd}P}{\mathrm{d}y^2}\end{aligned}
\]
Ut vero obtineamus aequationem ipsi 80 similem, statuamus $P=(1-y)^\mu P'$, unde
\[
\begin{aligned}[c]
\frac{\mathrm{d}P}{\mathrm{d}y}
&=-\mu(1-y)^{\mu-1}P'+(1-y)^\mu\frac{\mathrm{d}P'}{\mathrm{d}y}\\
\frac{\mathrm{dd}P}{\mathrm{d}y^2}
&=(\mu\mu-\mu)(1-y)^{\mu-2}P'-2\mu(1-y)^{\mu-1}\frac{\mathrm{d}P'}{\mathrm{d}y}
\\&\qquad+(1-y)^\mu\frac{\mathrm{dd}P'}{\mathrm{d}y^2}
\end{aligned}
\]
Quibus substitutis fit post divisionem per $(1-y)^\mu$
\[
\begin{aligned}[c]
&0=P'\{\alpha\beta-\mu(\gamma+(\alpha+\beta-1-\gamma)y+y(\mu\mu-\mu)\}\\
&+\frac{\mathrm{d}P'}{\mathrm{d}y}\{\gamma+(\alpha+\beta-1-\gamma)y-2\mu y\}(1-y)\\
&+\frac{\mathrm{dd}P'}{\mathrm{d}y^2}\{y-yy\}(1-y)
\end{aligned}
\]
Determinemus $\mu$ ita, ut multiplicator ipsius $P'$ per $1-y$ divisibilis evadat,
quod fiet statuendo vel $\mu=\alpha$ vel $\mu=\beta$. Valor prior mutat aequationem praecedentem in hanc
\[
\begin{aligned}&0=\alpha(\beta-\gamma)P'+(\gamma-(\gamma+\alpha+1-\beta)y)\frac{\mathrm{d}P'}{\mathrm{d}y}\\&+(y-yy)\frac{\mathrm{dd}P'}{\mathrm{d}y^2}\end{aligned}
\]
sive
\[
\begin{aligned}&0=\alpha(\gamma-\beta)P'-(\gamma-(\gamma-\beta+\alpha+1)y)\frac{\mathrm{d}P'}{\mathrm{d}y}\\&-(y-yy)\frac{\mathrm{dd}P'}{\mathrm{d}y^2}\end{aligned}
\]
cui ita satisfaciendum est, ut fiat pro $y=0$, $P'=1$ atque
$\frac{\mathrm{d}P'}{\mathrm{d}y}=\frac{\alpha(\gamma-\beta)}{\gamma}$. Hinc
autem deducitur $P'=F(\alpha,\gamma-\beta,\gamma,y)$ adeoque habetur
\[
\begin{aligned}&[91]\quad F(\alpha,\beta,\gamma,x)=(1-y)^\alpha F(\alpha,\gamma-\beta,\gamma,y)\\&=(1-x)^{-\alpha}F(\alpha,\gamma-\beta,\gamma,-\tfrac{x}{1-x})\end{aligned}
\]
Si pro $\mu$ valorem alterum $\beta$ adoptavissemus, prodiisset prorsus simili modo
\[
[92]\quad F(\alpha,\beta,\gamma,x)=(1-x)^{-\beta}F(\beta,\gamma-\alpha,\gamma,-\tfrac{x}{1-x})
\]
quae formula quoque e praecedenti per solam permutationem elementorum $\alpha$, $\beta$
sponte sequitur. Adiumento formulae modo inventae valores serierum nostrarum
pro valoribus negativis elementi quarti semper ad valores similium serierum pro
valoribus positivis elementi quarti interque $0$ et $1$ sitis reducitur, quum fiat
\[
F(\alpha,\beta,\gamma,-x)=(1+x)^{-\alpha}F(\alpha,\gamma-\beta,\gamma,\tfrac{x}{1+x})
\]
\SourceLeaf{228}
\SixthArticle{48}

Operae pretium erit ostendere, quomodo adiumento transformationum
82,91 omnes formulae in art. 5 collectae perfacile e solo theoremate binomiali
deduci possint. Protinus enim inde sequuntur formulae I---IV. Formulae VI---IX hinc sponte sequuntur, si \(e^x\) spectatur tamquam limes potestatis
\(\left(1+\frac{x}{i}\right)^i\) vel \(\left(1-\frac{x}{i}\right)^{-i}\), atque \(\log x\) tamquam limes ipsius \(i(x^{\frac{1}{i}}-1)\), crescente \(i\)
in infinitum. Porro ex
\[
\begin{aligned}[c]
\cos n\varphi+\BareRoot{-1}.\sin n\varphi&=(\cos\varphi+\BareRoot{-1}.\sin\varphi)^n\\
\cos n\varphi-\BareRoot{-1}.\sin n\varphi&=(\cos\varphi-\BareRoot{-1}.\sin\varphi)^n
\end{aligned}
\]
sequitur per subtractionem et additionem formula XVIII et XXII atque hinc
per formulam 82 statim XIX et XXIII; hinc rursus per formulam 91 deducuntur XVI, XVII, XX et XXI. Statuendo \(t\) pro \(nt\) atque \(n\) infinitum ex XVI et
XX sequuntur XI et XII; statuendo vero \(n\) infinite parvum ex XVI---XVIII
sequuntur XIII---XV.

\SixthArticle{49}

Ex substitutione \(x=\frac{1}{y}\) simili modo evadit
\[
\begin{aligned}&0=\alpha\beta P-(\alpha+\beta-1-(\gamma-2)y)y\frac{\mathrm{d}P}{\mathrm{d}y}\\&+(yy-y^3)\frac{\mathrm{d}\mathrm{d}P}{\mathrm{d}y^2}\end{aligned}
\]
Statuendo dein \(P=y^\mu P'\), fit
\[
\begin{aligned}[c]
&\mathrm{I.}\qquad 0=P'(\alpha\beta-\mu(\alpha+\beta-1)\\&\qquad+\mu(\gamma-2)y+(\mu\mu-\mu)(1-y))\\
&-\frac{\mathrm{d}P'}{\mathrm{d}y}(\alpha+\beta-1-(\gamma-2)y-2\mu(1-y))y\\
&+(yy-y^3)\frac{\mathrm{d}\mathrm{d}P'}{\mathrm{d}y^2}
\end{aligned}
\]
Ut multiplicator ipsius \(P'\) per \(y\) divisibilis evadat, statui debet vel \(\mu=\alpha\) vel
\(\mu=\beta\); valor prior producit
\[
\begin{aligned}&\mathrm{II.}\qquad
0=P'\alpha(\gamma-\alpha-1)-\frac{\mathrm{d}P'}{\mathrm{d}y}(\beta-\alpha-1-(\gamma-2\alpha-2)y)\\&+(y-yy)\frac{\mathrm{d}\mathrm{d}P'}{\mathrm{d}y^2}\end{aligned}
\]
cuius integrale \textit{particulare} fit
\[
P'=F(\alpha,\alpha+1-\gamma,\alpha+1-\beta,y)
\]
\par\SourceLeaf{229}
Satisfit igitur aequationi I per integrale particulare
\[
P=y^\alpha F(\alpha,\alpha+1-\gamma,\alpha+1-\beta,y).
\]
et perinde valor alter \(\mu=\beta\) suppeditat alterum integrale particulare
\[
P=y^\beta F(\beta,\beta+1-\gamma,\beta+1-\alpha,y)
\]
unde integrale completum habetur
\[
\begin{aligned}&P=Ay^\alpha F(\alpha,\alpha+1-\gamma,\alpha+1-\beta,y)\\&+By^\beta F(\beta,\beta+1-\gamma,\beta+1-\alpha,y)\end{aligned}
\]
designantibus \(A,B\) constantes, quae vero non sunt arbitrariae sed penitus determinatae, quum \(P\) non sit integrale completum aequationis 80, sed particulare tantum. At ne determinatio valorum constantium \(A,B\) in ambages inutiles nos deducat, eandem aequationem alio modo per ea, quae iam evolvimus, deducemus.

Statuendo in aequatione 91 \(-\frac{x}{1-x}=1-z\), mutandoque in aequat. 86
\(\beta\) in \(\gamma-\beta\), \(\gamma\) in \(\alpha+1-\beta\), \(x\) in \(z\), colligetur
\[
\begin{aligned}[c]
&(1-x)^\alpha F(\alpha,\beta,\gamma,x)\\&=
\frac{\Pi(\gamma-1)\Pi(\beta-\alpha-1)}{\Pi(\gamma-\alpha-1)\Pi(\beta-1)}F(\alpha,\gamma-\beta,\alpha+1-\beta,z)\\
&+\frac{\Pi(\gamma-1)\Pi(\alpha-\beta-1)}{\Pi(\alpha-1)\Pi(\gamma-\beta-1)}z^{\beta-\alpha}F(\beta,\gamma-\alpha,\beta+1-\alpha,z)
\end{aligned}
\]
Sed per aequatt. 91, 92
\[
\begin{aligned}[c]
&F(\alpha,\gamma-\beta,\alpha+1-\beta,z)\\&=(1-z)^{-\alpha}F\left(\alpha,\alpha+1-\gamma,\alpha+1-\beta,-\frac{z}{1-z}\right)\\[5pt]
&F(\beta,\gamma-\alpha,\beta+1-\alpha,z)\\&=(1-z)^{-\beta}F\left(\beta,\beta+1-\gamma,\beta+1-\alpha,-\frac{z}{1-z}\right)
\end{aligned}
\]
His substitutis, nec non \(z=\frac{1}{1-x}\), \(1-z=-\frac{x}{1-x}\), \(-\frac{z}{1-z}=\frac{1}{x}\), habemus
\[
\begin{aligned}[c]
&[93]\qquad F(\alpha,\beta,\gamma,x)\\&=
\frac{\Pi(\gamma-1)\Pi(\beta-\alpha-1)}{\Pi(\gamma-\alpha-1)\Pi(\beta-1)}(-x)^{-\alpha}F\left(\alpha,\alpha+1-\gamma,\alpha+1-\beta,\frac{1}{x}\right)\\
&+\frac{\Pi(\gamma-1)\Pi(\alpha-\beta-1)}{\Pi(\alpha-1)\Pi(\gamma-\beta-1)}(-x)^{-\beta}F\left(\beta,\beta+1-\gamma,\beta+1-\alpha,\frac{1}{x}\right)
\end{aligned}
\]
quae convenit cum aequatione supra inventa, si statuatur
\[
\begin{aligned}[c]
A&=\frac{\Pi(\gamma-1)\Pi(\beta-\alpha-1)}{\Pi(\gamma-\alpha-1)\Pi(\beta-1)}(-1)^\alpha\\
B&=\frac{\Pi(\gamma-1)\Pi(\alpha-\beta-1)}{\Pi(\alpha-1)\Pi(\gamma-\beta-1)}(-1)^\beta
\end{aligned}
\]
\SourceLeaf{230}
ubi notandum, esse
\[
\begin{array}{rcl}
(-1)^\alpha&=&\cos\alpha k\pi+\BareRoot{-1}\,.\sin\alpha k\pi\\
(-1)^\beta&=&\cos\beta k\pi+\BareRoot{-1}\,.\sin\beta k\pi
\end{array}
\]
designante \(k\) indefinite integrum imparem quemcunque.

\SixthArticle{50}
Per aequationem 93 valor functionis nostrae pro valoribus elementi quarti unitate maioribus ad casum eum reducitur, ubi elementum quartum unitate minus est. Simul patet, valoribus elementi quarti \textit{negativis}, unitate maioribus, semper unum valorem realem functionis \(F\) respondere, positivis vero tunc tantum respondere posse valorem realem functionis, si \(\alpha\) et \(\beta\) sint vel integri vel fractiones rationales, quarum denominatores sint impares; in casibus reliquis \(F(\alpha,\beta,\gamma,x)\) pro valore positivo ipsius \(x\) unitate maiore tantummodo valores imaginarios admittit.

\SixthArticle{51}
Relationes inter plures functiones \(F\) hactenus evolutae omnes lineares fuerunt: adiicimus aliam diversi generis. Sit
\[
\begin{array}{rcl}
P&=&F(\alpha,\beta,\gamma,x)\\
Q&=&x^{1-\gamma}F(\alpha+1-\gamma,\beta+1-\gamma,2-\gamma,x)\\
R&=&F(\alpha,\beta,\alpha+\beta+1-\gamma,1-x)
\end{array}
\]
ita ut \(P,\ Q.\ R\) sint tria integralia particularia aequationis 80, sive sit
\[
\begin{aligned}
\mathrm{I.}\quad&0=\alpha\beta P-(\gamma-(\alpha+\beta+1)x)\dfrac{\mathrm dP}{\mathrm dx}\\&\qquad-(x-xx)\dfrac{\mathrm{dd}P}{\mathrm dx^2}\\[4pt]
\mathrm{II.}\quad&0=\alpha\beta Q-(\gamma-(\alpha+\beta+1)x)\dfrac{\mathrm dQ}{\mathrm dx}\\&\qquad-(x-xx)\dfrac{\mathrm{dd}Q}{\mathrm dx^2}\\[4pt]
\mathrm{III.}\quad&0=\alpha\beta R-(\gamma-(\alpha+\beta+1)x)\dfrac{\mathrm dR}{\mathrm dx}\\&\qquad-(x-xx)\dfrac{\mathrm{dd}R}{\mathrm dx^2}
\end{aligned}
\]
Multiplicando aequationem primam per \(Q\), secundam per \(P\), fit subtrahendo
\[
\begin{aligned}&0=(\gamma-(\alpha+\beta+1)x)\frac{Q\mathrm dP-P\mathrm dQ}{\mathrm dx}\\&+(x-xx)\frac{Q\mathrm{dd}P-P\mathrm{dd}Q}{\mathrm dx^2}\end{aligned}
\]
Haec vero aequatio per \(x^{\gamma-1}(1-x)^{\alpha+\beta-\gamma}\) multiplicata integrabilis evadit atque suppeditat
\SourceLeaf{231}
\[
[94]\qquad A=x^\gamma(1-x)^{\alpha+\beta+1-\gamma}\frac{Q\mathrm dP-P\mathrm dQ}{\mathrm dx}
\]
Prorsus simili modo habetur
\[
\begin{array}{rl}
[95]&B=x^\gamma(1-x)^{\alpha+\beta+1-\gamma}\dfrac{R\mathrm dQ-Q\mathrm dR}{\mathrm dx}\\[11pt]
[96]&C=x^\gamma(1-x)^{\alpha+\beta+1-\gamma}\dfrac{R\mathrm dP-P\mathrm dR}{\mathrm dx}
\end{array}
\]
Constantes \(A,\ B,\ C\) facile determinantur per methodum sequentem.

Pro \(x=0\), fit \(P=1\); porro \(x^\gamma Q=xF(\alpha+1-\gamma,\beta+1-\gamma,2-\gamma,x)\) fit \(=0\) pro \(x=0\); ipsius differentiale autem per \(\mathrm dx\) divisum, puta \(\gamma x^{\gamma-1}Q+x^\gamma\frac{\mathrm dQ}{\mathrm dx}\) fit \(=1\); hinc colligitur \(\frac{x^\gamma\mathrm dQ}{\mathrm dx}=1-\gamma\) pro \(x=0\), adeoque
\[
A=\gamma-1
\]
Ut vero etiam \(B\) et \(C\) determinemus, resumamus aequationem
\[
R=f(\alpha,\beta,\gamma)P+f(\alpha+1-\gamma,\beta+1-\gamma,2-\gamma)Q
\]
quae differentiata dat
\[
\frac{\mathrm dR}{\mathrm dx}=f(\alpha,\beta,\gamma)\frac{\mathrm dP}{\mathrm dx}+f(\alpha+1-\gamma,\beta+1-\gamma,2-\gamma)\frac{\mathrm dQ}{\mathrm dx}
\]
Multiplicando primam per \(\frac{\mathrm dQ}{\mathrm dx}\), secundam per \(Q\), fit subtrahendo
\[
\frac{Q\mathrm dR-R\mathrm dQ}{\mathrm dx}=f(\alpha,\beta,\gamma)\frac{Q\mathrm dP-P\mathrm dQ}{\mathrm dx}\quad\mbox{adeoque}
\]
\[
B=(1-\gamma)f(\alpha,\beta,\gamma)=\frac{\Pi(\alpha+\beta-\gamma)\Pi(1-\gamma)}{\Pi(\alpha-\gamma)\Pi(\beta-\gamma)}
\]
Similiter multiplicata aequatione priore per \(\frac{\mathrm dP}{\mathrm dx}\), posteriore per \(P\), subtractio dat
\[
\frac{R\mathrm dP-P\mathrm dR}{\mathrm dx}=f(\alpha+1-\gamma,\beta+1-\gamma,2-\gamma)\frac{Q\mathrm dP-P\mathrm dQ}{\mathrm dx}
\]
adeoque
\[
\begin{aligned}&C=(\gamma-1)f(\alpha+1-\gamma,\beta+1-\gamma,2-\gamma)\\&=\frac{\Pi(\alpha+\beta-\gamma)\Pi(\gamma-1)}{\Pi(\alpha-1)\Pi(\beta-1)}\end{aligned}
\]
Si magis placet, hae tres aequationes etiam ita exhiberi possunt, ut functiones derivatae \(\frac{\mathrm dP}{\mathrm dx},\ \frac{\mathrm dQ}{\mathrm dx},\ \frac{\mathrm dR}{\mathrm dx}\) per functiones finitas exprimantur; ita e. g. fit formula 96,
